\documentclass[10pt,a4paper,openany,twoside,prof]{book}

\usepackage{layout_cours}

\setlist{leftmargin=5mm}

\begin{document}
	\chapnumber{1} % numéro du chapitre
	\titre{Ensembles de nombres} % titre du chapitre
	\classe{Seconde} % classe

\section*{Corrections (Indice Bordas 2019)}
\begin{bookexo}{115 p.69}\end{bookexo}
$\sqrt{11}\approx 3,316624$

On prend donc les 5 premiers chiffre pour obtenir $a=3.31662$.

On déduit $b=3.31663$.

On a l'encadrement:

$$a\le \sqrt{11} < b$$
$$3.31662\le \sqrt{11} < 3.31663$$

\begin{bookexo}{116 p.69}\end{bookexo}
$\frac{113}{31}\approx 3,64516129$

$3,6451612\le\frac{113}{31}<3,6451613$

\bigskip
$\frac{1}{17} \approx 0,06$

$0\le\frac{1}{17}<0.1$
\bigskip

\begin{bookexo}{120 p.69}\end{bookexo}
\begin{enumerate}[label=\alph*)]
	\item $\frac{\pi}{4}\in\R$ car $\pi\notin\Q$
	\item $10^{-4}\in\D$,\quad $10^{-4}\in\Q$ et $10^{-4}\in\R$ car $10^{-4}=0.0001$ (et $10^{-4}\notin\Z$ car il n'est pas entier)
	\item $-14.6\in\D$,\quad $-14.6\in\Q$ et $-14.6\in\R$ car il est déjà sous forme d'un nombre décimal (et $-14.6\in\Z$ car il n'est pas entier)
	\item $\frac{197}{2}\in\D$,\quad $\frac{197}{2}\in\Q$ et $\frac{197}{2}\in\R$ car $\frac{197}{2} = 98.5$ (et $98.5\notin\Z$ car il n'est pas entier)
\end{enumerate}
\end{document}
